\documentclass[11pt]{article}
\usepackage[a4paper,margin=18mm]{geometry}
\usepackage[no-math]{fontspec}
\setmainfont{Latin Modern Roman}
\usepackage{amsmath,amssymb,amsthm,xfrac,xspace,hyperref,graphicx,comment}
\excludecomment{DefTralics}
\providecommand{\dpl}{\displaystyle}
\providecommand{\backmatter}{}
\providecommand{\loccit}{\emph{loc. cit.}}

\providecommand{\bibliofont}{\small}
\newcommand{\ie}{i.e.,\xspace}
\newcommand{\eg}{e.g.,\xspace}
\newcommand{\etc}{etc.\xspace}
\newcommand{\cf}{cf.\xspace}
\newcommand{\resp}{resp.\xspace}
\newcommand{\smaller}{\small}
\newcommand{\Subsection}{\subsection}
\newcommand{\Subsubsection}{\subsubsection}
\newcommand{\skpt}{\smallskip}
\emergencystretch=3em
\setcounter{section}{1}\numberwithin{equation}{section}
% Begin literal retained input author-preamble.tex
\usepackage[scr=rsfs,cal=euler]{mathalfa}
\newenvironment{enumeratea}
{\bgroup\def\theenumi{\alph{enumi}}\begin{enumerate}}
{\end{enumerate}\egroup}

\let\epsilon\varepsilon
\newcommand\mto{\mathchoice{\longmapsto}{\mapsto}{\mapsto}{\mapsto}}
\newcommand\quand{\quad\text{and}\quad}
\newcommand\qquand{\qquad\text{and}\qquad}
\newcommand{\spfrac}[2]{\sfrac{#1}{(#2)}}

\usepackage{tikz}
\usetikzlibrary{cd}
\tikzcdset{arrow style=Latin Modern}
\newcommand{\dradpl}{{\hspace*{-2.5mm}\begin{tikzcd}[ampersand replacement=\&,column sep = scriptsize]\mbox{}\ar[r,dashed]\&\mbox{}\end{tikzcd}\hspace*{-2.5mm}}}
\newcommand{\dratst}{{\hspace*{-2.5mm}\begin{tikzcd}[ampersand replacement=\&,column sep = small]\mbox{}\ar[r,dashed]\&\mbox{}\end{tikzcd}\hspace*{-2mm}}}
\let\drascst\dratst
\newcommand{\dra}{\mathchoice{\dradpl}{\dratst}{\drascst}{\drascst}}


% Literal retained local macro stmaryrd



\usepackage[all,cmtip]{xy}
\xyoption{matrix}\let\labelstyle\textstyle

\usepackage{mathtools}

\newtheorem{mytheorem}{Theorem}
\renewcommand*{\themytheorem}{\Alph{mytheorem}}
\newtheorem{lemma}{Lemma}[section]
\newtheorem{corollary}[lemma]{Corollary}
\newtheorem{theorem}[lemma]{Theorem}
\newtheorem{proposition}[lemma]{Proposition}
\newtheorem{question}[lemma]{Question}

\newtheorem{maintheorem}{Theorem}
\renewcommand\themaintheorem{\Alph{maintheorem}}
\newtheorem{maincorollary}[maintheorem]{Corollary}
\newtheorem{mainprop}[maintheorem]{Proposition}

\theoremstyle{definition}
\newtheorem{definition}[lemma]{Definition}
\newtheorem{notation}[lemma]{Notation}

\newtheorem{remark}[lemma]{Remark}
\newtheorem{example}[lemma]{Example}

\newcommand\iso{\stackrel{\simeq}{\longrightarrow}}

\newcommand\A{{\mathbb A}}\let\AA\A
\newcommand\C{{\mathbb C}}\let\CC\C
\renewcommand\k{{\mathbf k}}
\newcommand\p{{\mathbb P}}\let\PP\p
\newcommand\Z{{\mathbb Z}}\let\ZZ\Z
\newcommand\FF{{\mathbb F}}
\newcommand\GG{{\mathbb G}}
\newcommand\NN{{\mathbb N}}
\newcommand\QQ{{\mathbb Q}}

\newcommand{\kk}{\mathbf{k}}

\let\rat\dra

\newcommand\mm{{\mathfrak m}}

\DeclareMathOperator{\Aut}{Aut}
\DeclareMathOperator{\car}{char}
\DeclareMathOperator{\sing}{sing}
\DeclareMathOperator{\reg}{reg}
\DeclareMathOperator{\Pic}{Pic}
\DeclareMathOperator{\Spec}{Spec}
\DeclareMathOperator{\Var}{Var}
\newcommand{\et}{\textup{ét}}

\newcommand{\set}[2]{\{#1 \mid #2\}}
\newcommand{\Bigset}[2]{\Bigl\{#1 \Bigm| #2\Bigr\}}

\hyphenation{co-cycle}

% End literal retained input author-preamble.tex

\begin{document}
\begin{center}{\Large Stable item 2248}\end{center}
\paragraph{Source and attribution.}
Jérémy Blanc, Pierre-Marie Poloni, Immanuel Van Santen, \emph{Complements of hypersurfaces in projective spaces}, Journal de l’École polytechnique — Mathématiques 11 (2024), publisher version, DOI 10.5802/jep.264. \href{https://jep.centre-mersenne.org/articles/10.5802/jep.264/}{Publisher article}. Authors' text: \href{https://creativecommons.org/licenses/by/4.0/}{CC BY 4.0}.
\paragraph{Editorial problem boundary.}
Prove only TheoremA below over the exact arbitrary field and integers d,n at least3, with characteristic not dividing d and the excluded pair d equals3,n equals3. Construct two nonisomorphic irreducible degree-d projective hypersurfaces with isomorphic complements. The complete new graded affine lifting/descent, explicit Danielewski-cylinder formulas, planar singular-curve distinction and exceptional-divisor/multiplicity induction are retained. No algebraic closedness is added. Low-degree positive results, Euler-characteristic questions and later cubic-surface classification are unselected and omitted.
\paragraph{Difficulty (editorial): H3.}
An affine cylinder isomorphism does not automatically descend to an isomorphism of projective complements, and it says nothing about whether the removed projective hypersurfaces are isomorphic. The authors construct inverse polynomial formulas with controlled homogeneous degree classes, then distinguish the hypersurfaces using singular local rings and maximal-multiplicity exceptional divisors through higher-dimensional cones. These new descent and geometric-invariant arguments are necessary before the final case split, beyond the classical cyclic-quotient or blow-up foundations alone.
\paragraph{Editorial transcription note.}
Original author language, mathematics and ordering are preserved. Publisher front matter is replaced by this independent article wrapper. Bibliography is recovered from matching cached publisher HTML, including its TeX formulas. No new proof or mathematical repair is supplied. Surrounding original results are retained for conservative context and are not extra selected corpus claims.
\clearpage

% Begin literal retained input original-primary-statement.tex
\begin{maintheorem}\label{34and43}
Let $\kk$ be a field and let $d,n \ge 3$ be integers such that $d$ is not a multiple of $\mathrm{char}(\kk)$ and such that $(d,n)\neq (3,3)$.

There are two non-isomorphic irreducible hypersurfaces $H,H'\subseteq \p^n_\k$ of degree $d$, that have isomorphic complements.
\end{maintheorem}

% End literal retained input original-primary-statement.tex


% Begin literal retained input author-body.tex
\section{Lift to affine hypersurfaces}
Let $n\geq 1$. For each homogeneous polynomial $f\in \kk[x_0,\ldots,x_n]$ of degree $d\ge 1$ and each $\mu\in \k^*$, we denote by $X_{f,\mu}\subseteq \A^{n+1}$ the hypersurface given by
\[
X_{f,\mu}=\Spec(\k[x_0,\ldots,x_n]/(f-\mu)) \subseteq \AA^{n+1},
\]
and obtain a canonical finite morphism
\begin{align*}
\pi_{f,\mu}\colon X_{f, \mu}&\to \p^n_f\\
(x_0,\ldots,x_n)&\mto [x_0:\cdots:x_n]
\end{align*}
of degree $d$, which is an étale covering if $\car(\kk)$ does not divide $d$.

If $\kk$ is algebraically closed, all the varieties $X_{f,\mu}$ are isomorphic to $X_{f,1}$ by homotheties.

If $\k$ is equal to the field of complex numbers $\CC$, then $\pi_{f,\mu}$
is the universal abelian covering, as the abelianization of the
fundamental group of $\p^n_f$ with respect to the Euclidean topology is equal to $\ZZ / d \ZZ$ \cite[Prop.\,2.3]{Li2005Lectures-on-topolo}.

We now prove that for each field $\kk$, the isomorphisms between complements of hypersurfaces lift to these coverings. This will be useful in the sequel, in order to study isomorphisms between varieties $\p^n_f$ and $\p^n_g$.

\begin{proposition}
\label{Prop.cyclic_covering}
Let $\kk$ be a field, let $n \geq 1$ and let $f,g\in \k[x_0,\ldots,x_n]$ be irreducible homogeneous polynomials
of degree $d \ge 1$. Let $\Phi_0,\ldots,\Phi_n\in \k[x_0,\ldots,x_n]$ be homogeneous polynomials of degree $\ell\ge 1$.
Then, the following statements are equivalent:
\begin{enumerate}
\item\label{cyclic_covering1} The rational map
\begin{align*}
\Phi\colon \p^n_f & \iso  \p^n_g, \\
{[x_0:\cdots:x_n]} & \mto  [\Phi_0(x_0,\ldots,x_n):\cdots: \Phi_n(x_0,\ldots,x_n)]
\end{align*}
is an isomorphism and the gcd of $\Phi_0, \ldots, \Phi_n$ is a power of $f$.
\item\label{cyclic_covering2}There exists $\mu\in \k^*$ such that the following map is an isomorphism
\begin{align*}
\varphi\colon X_{f,1} & \iso X_{g,\mu}, \\
(x_0,\ldots,x_n) & \mto (\Phi_0(x_0,\ldots,x_n),\ldots, \Phi_n(x_0,\ldots,x_n)).
\end{align*}
\end{enumerate}
Moreover, if these statements hold, then $\ell$ is invertible modulo $d$.
\end{proposition}
\begin{proof}
``\eqref{cyclic_covering1} $\Rightarrow$ \eqref{cyclic_covering2}'':
There exists $r \geq 0$ and $\Phi_0', \ldots, \Phi_n'\in\kk[x_0, \ldots, x_n]$
without common factor such that
$\Phi_i = f^r \Phi_i'$ for all $i$. The domain $\textrm{dom}(\Phi)$ of the rational self-map $\Phi \colon \PP^n \rat \PP^n$
is thus equal to
$\PP^n \setminus V_{\PP^n}(\Phi'_0, \ldots, \Phi'_n)$ and contains $\PP^n_f$.
As $g$ is nowhere zero on $\p^n_g$, the polynomial $g(\Phi_0,\ldots,\Phi_n) = f^{dr} g(\Phi_0', \ldots, \Phi_n')$ is nowhere zero on $\p^n_f$ and is thus equal to $\mu f^t$ for some $\mu\in \k^*$ and some integer $t\ge 1$. We then obtain a morphism
\[
\varphi\colon X_{f,1} \to X_{g,\mu}, \quad
(x_0,\ldots,x_n)\mto (\Phi_0(x_0,\ldots,x_n),\ldots, \Phi_n(x_0,\ldots,x_n)).
\]
It remains to see that $\varphi$ is an isomorphism.

The inverse of $\Phi$ is given by
\[
\Psi\colon \p^n_g\iso \p^n_f, \quad [x_0:\cdots:x_n]\mto [\Psi_0(x_0,\ldots,x_n):\cdots: \Psi_n(x_0,\ldots,x_n)],
\]
where $\Psi_0,\ldots,\Psi_n\in \k[x_0,\ldots,x_n]$ are homogeneous polynomials of the same degree $\ell'\ge 1$, without common factor.

As $\Psi\circ \Phi=\mathrm{id}_{\p^n_f}$ and $\Phi\circ \Psi=\mathrm{id}_{\p^n_g}$, there exist homogeneous polynomials $A,B\in \k[x_0,\ldots,x_n]$ of degree $\ell\ell'-1$ such that
\[\Psi_i(\Phi_0,\ldots,\Phi_n)= Ax_i \quand\Phi_i(\Psi_0,\ldots,\Psi_n)= Bx_i \qquad\text{for all }i\in \{0,\ldots,n\}.\]
Hence, the zero locus $V_{\PP^n_f}(A)$ of $A$ in $\PP^n_f$ satisfies the following: $\Psi_i$ vanishes on $\Phi(V_{\PP^n_f}(A))$ for all
$i \in \{0, \ldots, n\}$
and $\Phi(V_{\PP^n_f}(A))$ lies in the domain of $\Psi$; \ie $V_{\PP^n_f}(A)$
is empty. Analogously, one can see that $V_{\PP^n_g}(B)$ is empty.
We get $A=\lambda f^s$ and $B=\lambda' g^{s}$ where $\lambda,\lambda'\in \k^*$ and where $s\ge 1$ is such that $ds=\ell\ell'-1$. In particular, $\ell$ is invertible modulo $d$.
We replace $\Psi_i$ with $\Psi_i/\lambda$ for each $i\in \{0,\ldots,n\}$ and reduce to the case where $\lambda=1$.

Consider now the morphism
\[
\psi\colon X_{g,\mu} \to \A^{n+1}, \quad
(x_0,\ldots,x_n)\mto (\Psi_0(x_0,\ldots,x_n),\ldots, \Psi_n(x_0,\ldots,x_n)).
\]
Using that $\Psi_i(\Phi_0,\ldots,\Phi_n)= f^sx_i$ for each $i\in \{0,\ldots,n\}$, we obtain that
$\psi\circ \varphi \colon X_{f, 1} \to \AA^{n+1}$ is the natural closed embedding.
Hence, $(\psi \circ \varphi)^\ast = \varphi^\ast \circ \psi^\ast \colon \kk[\AA^{n+1}] \to \kk[X_{f, 1}]$ is surjective and $\varphi$ is dominant, \ie $\varphi^\ast \colon \kk[X_{g, \mu}] \to \kk[X_{f, 1}]$ is injective.
This implies that
$\varphi^\ast$ and thus $\varphi$ is an isomorphism.

``\eqref{cyclic_covering2} $\Rightarrow$ \eqref{cyclic_covering1}'': As $\Phi_0, \ldots, \Phi_{n}$ are homogeneous of degree $\ell$, we may define a rational self-map $\Phi\colon [x_0:\cdots:x_n]\mto [\Phi_0(x_0,\ldots,x_n):\cdots: \Phi_n(x_0,\ldots,x_n)]$ of $\p^n$. We now prove that $\Phi$ restricts to an isomorphism $\p^n_f\iso \p^n_g$ and that the gcd
of $\Phi_0, \ldots, \Phi_n$ is a power of $f$.

By assumption $g(\Phi_0, \ldots, \Phi_n) = (f-1) p + \mu$ for some polynomial $p \in \kk[x_0, \ldots, x_n]$
of degree $r = d (\ell-1) \geq 0$.
Write $p = \sum_{i=0}^r p_i$ where $p_i$ is homogeneous of degree $i$.
Since $g(\Phi_0, \ldots, \Phi_n)$ is homogeneous, it~follows that $p_0 = \mu$ and $p_{di} = f p_{d(i-1)}$ for all
$1 \leq i \leq \ell-1$. This implies that
\[
g(\Phi_0, \ldots, \Phi_n) = \mu f^{\ell}.
\]
Since $g$ is homogeneous
and $f$ is irreducible the gcd of $\Phi_0, \ldots, \Phi_n$ is a power of $f$.
Moreover, $\Phi$ restricts to a morphism $\Phi \colon \p^n_f\to \p^n_g$.

Using that $\pi_{g, \mu} \circ \varphi = \Phi \circ \pi_{f, 1}$, we get the following commutative diagram
\[
\xymatrix{
\kk(X_{g, \mu}) \ar[r]^-{\varphi^\ast} & \kk(X_{f, 1}) \\
\kk(\PP^n_g) \ar[u]^-{\pi_{g, \mu}^\ast} \ar[r]^-{\Phi^\ast} & \kk(\PP^n_f) \ar[u]_-{\pi_{f, 1}^\ast}
}
\]
and that $\Phi \colon \p^n_f\to \p^n_g$ is a quasi-finite surjection.
Note that $\pi_{f, 1}^\ast$ and $\pi_{g, \mu}^\ast$ are field extensions of degree $d$.
Since $\varphi^\ast$ is an isomorphism, $\Phi^\ast$ is an isomorphism as well, \ie $\Phi \colon \PP^n_f \to \PP^n_g$
is birational. Zariski's Main Theorem \cite[Cor.\,4.4.9]{Gr1961Elements-de-geomet-III}
gives that $\Phi \colon \PP^n_f \to \PP^n_g$ is an isomorphism.
\end{proof}

\begin{remark}\label{Rem.mu_equal_1}
If there exists an isomorphism $\Phi \colon \PP^n_f \to \PP^n_g$, there are homogeneous polynomials
$\Phi_0,\ldots,\Phi_n\in \k[x_0,\ldots,x_n]$ of degree $\ell \geq 1$ without common factor
such that
$\Phi([x_0: \cdots: x_n]) = [\Phi_0(x_0,\ldots,x_n):\cdots: \Phi_n(x_0,\ldots,x_n)]$.
Proposition~\ref{Prop.cyclic_covering} then gives the existence of $\mu\in \k^*$ such that
\begin{align*}
\varphi\colon X_{f,1} & \iso X_{g,\mu}, \\
(x_0,\ldots,x_n) & \mto (\Phi_0(x_0,\ldots,x_n),\ldots, \Phi_n(x_0,\ldots,x_n))
\end{align*}
is an isomorphism. If $\kk$ is moreover algebraically closed, then
we may multiply each~$\Phi_i$ with a root of $\mu^{-1}$ and assume that $\mu=1$. Hence, $X_{f,1}$ and $X_{g,1}$ are isomorphic.
\end{remark}

\begin{remark}\label{remark:Replacemodulo}
Let $\kk$ be a field, $n \geq 1$ be an integer and let $f,g\in \k[x_0,\ldots,x_n]$ be irreducible homogeneous polynomials of degree $d \ge 1$. Let $e \ge 1$ be an integer and let $\Phi_0,\ldots,\Phi_n\in \k[x_0,\ldots,x_n]$ be polynomials such that the degree of each of their monomials is congruent to $e$ modulo $d$. If
\[
X_{f, 1} \iso X_{g, 1}, \quad (x_0,\ldots,x_n)\mto (\Phi_0(x_0,\ldots,x_n),\ldots, \Phi_n(x_0,\ldots,x_n))
\]
is an isomorphism, then we may multiply the monomials of $\Phi_i$ with some powers of~$f$ and assume that $\Phi_0,\ldots,\Phi_n$ are homogeneous of the same degree $\ell \geq 1$, and then apply Proposition~\ref{Prop.cyclic_covering} to obtain an isomorphism $\p^n_f\iso \p^n_g$.
\end{remark}

\begin{lemma} \label{lemma:Replace_n_by_m}
Let $\kk$ be a field, $n \geq 1$ be an integer, $f,g\in \k[x_0,\ldots,x_n]$ be irreducible homogeneous polynomials of degree $d \ge 1$, and let $\Phi\colon \p^n_f \iso \p^n_g$ be an isomorphism given by homogeneous polynomials of degree $\ell$ such that their gcd is a power of $f$.
If~$\ell$ is congruent to $1$ modulo $d$, then $\p^{m}_f$ and $\p^{m}_g$ are isomorphic for each $m\ge n$.
\end{lemma}

\begin{proof}
We write $\Phi$ as $[x_0:\cdots:x_n] \mto [\Phi_0(x_0,\ldots,x_n):\cdots: \Phi_n(x_0,\ldots,x_n)]$, where $\Phi_0,\ldots,\Phi_n\in \k[x_0,\ldots,x_n]$ are homogeneous of degree $\ell$ and the $\mathrm{gcd}$ of $\Phi_0,\ldots,\Phi_ n$ is a power of $f$. By Proposition~\ref{Prop.cyclic_covering}, the morphism
\[
\varphi\colon X_{f,1} \iso X_{g,\mu}, \quad (x_0,\ldots,x_n) \mto (\Phi_0(x_0,\ldots,x_n),\ldots, \Phi_n(x_0,\ldots,x_n))
\]
is an isomorphism, where $X_{f,1}, X_{g,\mu}\subseteq \A^{n+1}$ are defined as before. If $\ell$ is congruent to~$1$ modulo~$d$, we can write $\ell=1+dr$ for some integer $r\ge 0$. For each $m\ge n$,
we~thus get an isomorphism
$X_{f,1}\times \A^{m-n}\iso X_{g,\mu}\times \A^{m-n}$ given by
\[
(x_0,\ldots,x_m) \mto (\varphi(x_0,\ldots,x_m), x_{n+1}f(x_0,\ldots,x_n)^r,\ldots, x_m f(x_0,\ldots,x_n)^r)).
\]
Applying again Proposition~\ref{Prop.cyclic_covering}, now in the converse direction, we get that the homogeneous polynomials $\Phi_0,\ldots,\Phi_n, x_{n+1} f^r,\ldots,x_m f^r$ of degree $\ell$ also define an isomorphism $\p^{m}_f\iso\p^{m}_g$.
\end{proof}
The next examples shows that $\ell\neq 1$ in Proposition~\ref{Prop.cyclic_covering}
is possible.

\begin{example}
Let $\k$ be a field and let $f=g=xyz+x^3+y^3\in \k[x,y,z]$, which is the equation of a nodal cubic curve $\Gamma\subset\p^2$. By blowing-up the singular point, then blowing-up successively points on the exceptional divisor created lying on the curve, we obtain a birational morphism $X\to\p^2$ with a strict transform $\tilde{\Gamma}\subseteq X$
of $\Gamma$ being a $(-1)$-curve. We then contract this curve and contract the exceptional divisors created, except the last one, to get an automorphism of $\p^2\setminus \Gamma$ that does not extend to $\p^2$ and has degree $8$. Explicitly, we define $\Phi_0,\Phi_1,\Phi_2\in \k[x,y,z]$, homogeneous of degree $8$, by
\begin{align*}
\Phi_0 & = (-x^4z + 2x^3y^2 - 2x^2yz^2 + 2xy^3z + y^5 - y^2z^3)\cdot f, \\
\Phi_1 & = (x^2+yz)\cdot f^2, \\
\Phi_2 &=
x^7 y -x^6 z^2+6 x^5 y^2z-x^4 y^4-3 x^4 yz^3 +9 x^3 y^3z^2 \\
&\hspace*{2.14cm} + x^2y^5 z-3 x^2 y^2z^4 - x y^7+4 x y^4z^3 +2 y^6z^2- y^3z^5,
\end{align*}
and calculate that these define an involution
\[
\p^2_f\iso \p^2_f, \quad [x:y:z]\mto [\Phi_0(x,y,z):\Phi_1(x,y,z):\Phi_2(x,y,z)].
\]
Here, the common degree is $8$, that is not congruent to $1$ modulo $3$, but invertible in $\ZZ/3\ZZ$, as Proposition~\ref{Prop.cyclic_covering} says.
\end{example}

\begin{lemma}
\label{Lem.Open_iso_to_Pn_g}
Let $\kk$ be a field, $n\ge 1$, $f \in \kk[x_0, \ldots, x_n]$ an irreducible homogeneous polynomial of degree $d \geq 1$
and let $U \subseteq \p^n$ be an open subvariety.
If there exists an isomorphism $\p^n_f \iso U$, then there exists an irreducible homogeneous polynomial $g \in \kk[x_0, \ldots, x_n]$ of degree $d$ with $U = \p^n_g$.
\end{lemma}

Before we start the proof we recall the following fact: If $X$ is an affine irreducible
normal variety over a field $\kk$ and if $U \subseteq X$ is a big open subset, \ie the codimension of $X \setminus U$ in $X$ is at least $2$,
then $\kk[U] = \kk[X]$ (see \cite[Th.\,11.5]{Ma1986Commutative-ring-t}). In particular, if $U$ is affine as well, then $U = X$.

\begin{proof}[Proof of Lemma~$\ref{Lem.Open_iso_to_Pn_g}$]
If $U$ is big open in $\PP^n$, then for any open affine subvariety $V \subseteq\nobreak \PP^n$,
the affine subvariety $U \cap V$ is big open in $V$ and hence $U \cap V = V$. This would imply that $U = \PP^n$
is affine, a contradiction since $n \geq 1$. Hence the union $Z$ of the $(n-\nobreak1)$-dimensional
irreducible components of $\PP^n \setminus U$ is non-empty. As~$\PP^n \setminus Z$ is an affine
big open subset of $U$ we get $\PP^n \setminus Z = U$.
Moreover, $\kk[\p^n_f]^\ast= \k^\ast$ yields $\kk[U]^\ast=\nobreak\k^\ast$, which implies that $Z$ is irreducible.
Hence, there exists an irreducible homogeneous polynomial $g \in \kk[x_0, \ldots, x_n]$ such that $Z$
is the zero locus of $g$, \ie $U =\nobreak \PP^n_g$. Finally, $\ZZ / \deg(g) \ZZ \simeq \Pic(\p^n_g) \simeq \Pic(\p^n_f) \simeq \ZZ / d \ZZ$ implies that $\deg(g) =\nobreak d$.
\end{proof}

\section{Non-isomorphic hypersurfaces having isomorphic complements}
\subsection{Explicit isomorphisms between complements of projective cones}

The next formulas are
inspired by isomorphisms between cylinders over Danielewski surfaces given in \cite{MoPo2021Isomorphisms-betwe}.

\begin{lemma}\label{Lemma-Iso}
Let $\kk$ be a field, let $s,m,n\geq0$ be integers and let $S=\k[x_0,x_1,\ldots,x_s]$ be a polynomial ring in $s+1$ variables over $\kk$.
Let $P,Q\in S[z]$ and denote by $H_P$ and~$H_Q$ the hypersurfaces of $\A^{s+3}_{\k}=\Spec(S[y,z])$ defined by the equations
\[
x_0^ny=P(z) \quad \text{and}
\quad x_0^ny=Q(z),
\]
respectively. Suppose that there exist $A,B\in S[z]$ such that
\begin{enumerate}
\item $z-A(B(z))$ and $z-B(A(z))$ both belong to the ideal $x_0^{m}S[z]$,
\item We have
\begin{align*}
Q\left(A(z)+x_0^{m}w\right)&\in x_0^nS[z,w]+P(z)S[z,w]\\
\tag*{and}
P\left(B(z)+x_0^{m}w\right)&\in x_0^nS[z,w]+Q(z)S[z,w],
\end{align*}
where $w$ is transcendental over $S[y,z]$.
\end{enumerate}
Then, the following are inverse isomorphisms.
\begin{align*}
\Phi\colon H_P\times\A^1_{\k} &\iso H_Q\times\A^1_{\k}\\
(x_0,\ldots,x_s,y,z,w) &\mto \Bigl(x_0,\ldots,x_s,\frac{Q\left(A(z)+x_0^{m}w\right)}{x_0^n},  A(z)+x_0^{m}w,\\
&\hspace*{5cm}\frac{z-B\left(A(z)+x_0^{m}w\right)}{x_0^{m}}\Bigr),\\
\Psi\colon H_Q \times\A^1_{\k} &\iso H_P\times\A^1_{\k}\\
(x_0,\ldots,x_s,y,z,w) &\mto \Bigl(x_0,\ldots,x_s,\frac{P\left(B(z)+x_0^{m}w\right)}{x_0^n}, B(z)+x_0^{m}w,\\
&\hspace*{5cm}\frac{z-A\left(B(z)+x_0^{m}w\right)}{x_0^{m}}\Bigr).
\end{align*}
\end{lemma}

\begin{proof}
The hypotheses on $A$ and $B$ imply that $\Phi$ and $\Psi$ are morphisms. Moreover, the compositions are given by
\begin{align*}
(\Psi\circ\Phi)(x_0,\ldots,x_s,y,z,w)&=(x_0,\ldots,x_s,P(z)/x_0^n,z,w),\\
(\Phi\circ\Psi)(x_0,\ldots,x_s,y,z,w)&=(x_0,\ldots,x_s,Q(z)/x_0^n,z,w).\qedhere
\end{align*}
\end{proof}

\begin{proposition}\label{Prop.Counterexample.dim3}Let $\kk$ be a field, $d\ge 4$ be integers and let $f,g\in\k[x,y,z]$ be the homogeneous polynomials of degree $d$ defined by \[f=x^{d-1}y+z^{d} \quand g=x^{d-1}y+z^{d}+dx^{d-2}z^2.\] Then, for each $r\ge 1$, the
open subvarieties $\p^{r+2}_f,\p^{r+2}_g \subseteq \p^{r+2}$ are isomorphic.
\end{proposition}

\begin{remark}
The open subvarieties $\p^{2}_f,\p^{2}_g \subseteq \p^{2}$ are not isomorphic when $\kk=\mathbb{C}$ is the field of complex numbers. This follows from Proposition~\ref{Prop.cyclic_covering} and from the fact that, for every $\mu\in\C^*$, the (Danielewski) hypersurfaces $X_{f,1}, X_{g,\mu}\subseteq\A^3_{\mathbb{C}}$ defined by the equations $f=1$ and $g=\mu$, respectively, are not isomorphic (see \cite[Th.\,9]{PoloniClassifications}).
\end{remark}

\begin{proof}
We define $s=0$, $x_0=x$, $S=\kk[x]$, $m=d$, $n=d-1$,
\begin{align*}
P(z)&=1-z^{d}, & Q(z) &=1-z^{d}-dx^{d-2}z^2, \\
A(z)&=z-x^{d-2}z^3, & B(z) &=z+x^{d-2}z^3,
\end{align*}
and prove that the hypotheses of Lemma~\ref{Lemma-Iso} are satisfied for the above polynomials in $S[z]$.

Firstly, we have $A(z)\equiv z\equiv B(z) \pmod{x^{d-2}}$. Thus, since $d\le 2(d-2)$, as $d\ge 4$, we get $x^{d-2}A(z)^3\equiv x^{d-2}z^3\equiv x^{d-2} B(z)^3 \pmod{x^d}$ and we obtain
\begin{align*}
z-A(B(z))&= z-B(z)+x^{d-2}B(z)^3\equiv 0\pmod{x^d}, \\
z-B(A(z))&= z-A(z)-x^{d-2}A(z)^3\equiv 0\pmod{x^d}.
\end{align*}

We then check that $P\left(B(z)+x^{d}w\right)\in x^{d-1}S[z,w]+Q(z)S[z,w]$:
\begin{align*}
P(B(z)+x^{d}w) &\equiv P(B(z)) \\
&\equiv 1 -(z+x^{d-2}z^3)^{d} \\
&\equiv 1-z^{d}-dx^{d-2}z^{d+2}\\
&\equiv (1-z^{d}-dx^{d-2}z^2)\cdot (1+dx^{d-2}z^2)\\
&\equiv Q(z) \cdot (1+dx^{d-2}z^2)\, \pmod{x^{d-1}}.
\end{align*}
Similarly, we have
$Q(A(z)+x^{d}w)\in x^{d-1}S[z,w]+P(z)S[z,w]$:
\begin{align*}
Q(A(z)+x^{d}w) &\equiv Q\left(A(z)\right)\\
& \equiv 1-(z-x^{d-2}z^3)^{d}-dx^{d-2}(z-x^{d-2}z^3)^2 \\
&\equiv 1 -z^{d}+dx^{d-2}z^{d+2}-dx^{d-2}z^2\\
&\equiv (1-z^{d})\cdot (1-dx^{d-2}z^2)\\
&\equiv P(z) \cdot (1-dx^{d-2}z^2)\, \pmod{x^{d-1}}.
\end{align*}
Noting that $f-1=x^{d-1}y-P(z)$ and $g-1=x^{d-1}y-Q(z)$, we may now apply Lemma~\ref{Lemma-Iso} to obtain an isomorphism
\begin{align*}
\varphi\colon\Spec(\k[x,y,z,w]/(f-1)) & \iso \Spec(\k[x,y,z,w]/(g-1)) \\
(x,y,z,w) & \mto \Bigl(x,\frac{Q(A(z)+x^{d}w)}{x^{d-1}},
A(z)+x^{d}w,\\
&\hspace*{5cm}\frac{z-B(A(z)+x^{d}w)}{x^{d}}\Bigr).
\end{align*}

Arguing as in Remark~\ref{remark:Replacemodulo}, we may then use this isomorphism to construct an isomorphism $\p^3_f\iso \p^3_g$. Indeed, the second component of $\varphi$ can be replaced by a polynomial, since
\[
Q\left(A(z)+x^{d}w\right)\equiv P\cdot (1-dx^{d-2}z^2)\pmod{x^{d-1}}
\]
and $P\equiv x^{d-1}y\pmod{f-1}$.
Moreover, the fourth component of $\varphi$ is a polynomial, since $z-A(B(z)) \equiv 0\pmod{x^d}$.
Doing this, we obtain an expression for $\varphi$ given by four polynomials. It remains to check that each monomial appearing in that polynomial expression of $\varphi$ has degree congruent to $1$ modulo $d$. Working modulo $d$ with the degrees, we observe that $A,B,P,Q$ are homogeneous of degree $1,1,0,0$, respectively, hence that the polynomials $z-B(A(z)+x^{d}w)$, $Q(A(z)+x^dw)$ and $P-x^{d-1}y$ are homogeneous of degree $1$, $0$ and $0$, respectively. Therefore, we can conclude that $\p^3_f\iso \p^3_g$ as in Remark~\ref{remark:Replacemodulo} and furthermore that $\p^{r+2}_f\iso \p^{r+2}_g$ for all $r\geq 1$ by~Lemma~\ref{lemma:Replace_n_by_m}.
\end{proof}

\begin{proposition}
\label{Prop.Counterexample.dim4}
Let $\kk$ be a field, $d\ge 3$ be an integer and let $f,g\in\k[x_0,x_1,y,z]$ be the homogeneous polynomials of degree $d$ defined by
\[
f=x_0^{d-1}y+z^{d} \quand
g=x_0^{d-1}y+z^{d}+dx_0^{d-2}x_1^2. \]
Then, for each $r\ge 1$, the open subvarieties $\p^{r+3}_f,\p^{r+3}_g \subseteq \p^{r+3}$ are isomorphic.
\end{proposition}

\begin{remark}
We don't know, whether the open subvarieties $\p^{3}_f,\p^{3}_g \subseteq \p^{3}$ are isomorphic,
or even if the hypersurfaces $X_{f,1}, X_{g,1}\subseteq \A^3$ given by $f=1$ and $g=1$, respectively, are isomorphic or not.
\end{remark}

\begin{proof}[Proof of Proposition~$\ref{Prop.Counterexample.dim4}$]
We define $s=1$, $S=\kk[x_0,x_1]$, $m=d$, $n=d-1$, $\Delta=x_0^{d-2}x_1^2$,
\begin{align*}
P(z) &=1-z^{d}, & Q(z) &=1-z^{d}-d\Delta,\\
A(z) &=z(1-\Delta+\Delta^2), & B(z) &=z(1+\Delta),
\end{align*}
and prove that the hypotheses of Lemma~\ref{Lemma-Iso} are satisfied for the above polynomials in $S[z]$.

Firstly, we calculate $A(B(z))=B(A(z))=z(1-\Delta+\Delta^2)(1+\Delta)=z(1+\Delta^3)$.
As $d \geq 3$ we have $3d-6\ge d$, so $\Delta^3$ is divisible by $x_0^d$.
Hence, $A(B(z))=B(A(z))$ is congruent to $z$ modulo $x_0^d$.

We then check that $P(B(z)+x_0^{d}w)\in x_0^{d-1}S[z,w]+Q(z)S[z,w]$. As $d\ge 3$, $\Delta^2$ is divisible by $x_0^{d-1}$. Hence:
\begin{align*}
P(B(z)+x_0^{d}w) &\equiv P(B(z))\\
&\equiv 1-z^d(1+\Delta)^{d} \\
&\equiv 1- z^d(1+d\Delta) \\
&\equiv(1-z^{d}-d\Delta)\cdot (1+d\Delta)\\
&\equiv Q\cdot (1+d\Delta)\, \pmod{x_0^{d-1}}.
\end{align*}

Similarly, we have
$Q(A(z)+x_0^{d}w)\in x_0^{d-1}S[z,w]+P(z)S[z,w]$:
\begin{align*}
Q(A(z)+x_0^{d}w) &\equiv Q(A(z))\\
& \equiv 1-z^d(1-\Delta)^d-d\Delta\\
& \equiv 1-z^d(1-d\Delta)-d\Delta\\
& \equiv (1-z^{d})\cdot (1-d\Delta)\\
&\equiv P\cdot (1-d\Delta)\, \pmod{x_0^{d-1}}.
\end{align*}

Since $f-1=x_0^{d-1}y-P(z)$ and $g-1=x_0^{d-1}y-Q(z)$, we may now apply Lemma~\ref{Lemma-Iso} to obtain an isomorphism
\begin{align*}
\varphi\colon\Spec(\k[x_0,x_1,y,z,w]/(f-1)) & \iso \Spec(\k[x_0,x_1,y,z,w]/(g-1)) \\
(x_0,x_1,y,z,w) & \mto \Bigl(x_0,x_1,\frac{Q(A(z)+x_0^{d}w)}{x_0^{d-1}},A(z)+x_0^{d}w, \\
&\hspace*{4cm}\frac{z-B(A(z)+x_0^{d}w)}{x_0^{d}}\Bigr).
\end{align*}

We again proceed as in Remark~\ref{remark:Replacemodulo}.
The third component of $\varphi$ can be expressed by a polynomial, since
\[
Q(A(z)+x_0^{d}w)\equiv P\cdot (1- d x_0^{d-2}x_1^2)\pmod{x_0^{d-1}}
\]
and $P\equiv x_0^{d-1}y\pmod{f-1}$.
The last component of $\varphi$ is already a polynomial, as its numerator is a multiple of its denominator. It remains to check that each monomial appearing in the expression of $\varphi$ has degree congruent to $1$ modulo $d$. For this, we simply work with the degrees modulo $d$ and observe that $\Delta$ is homogeneous of degree $0$, so that $A,B,P,Q$ are homogeneous of degree $1,1,0,0$, respectively.
Hence, $Q(A(z)+x_0^dw)$, $A(z) + x_0^d w$ and $z-B(A(z)+x_0^{d}w)$ are homogeneous of degree
$0$, $1$ and $1$, respectively.
Therefore, we can conclude by Remark~\ref{remark:Replacemodulo} and
Lemma~\ref{lemma:Replace_n_by_m}.
\end{proof}

\subsection{Non-isomorphic hypersurfaces}
\label{Subsec.Counter_examples}
We now want to prove that the hypersurfaces of Propositions~\ref{Prop.Counterexample.dim3} and \ref{Prop.Counterexample.dim4} are not isomorphic (except when the characteristic divides $d$, in which case both are in fact equal).

Recall that the multiplicity $\mathrm{mult}_p(X)$ of a variety $X$ at a point $p$ is the multiplicity of the maximal ideal $\mathfrak{m}_{X,p}$ in the local ring $\mathcal{O}_{X,p}$.
If $X \subseteq \AA^n$ is a hypersurface given by a reduced polynomial
$f \in \kk[x_1, \ldots, x_n]$ of degree $d$, then $X$ has multiplicity
$d$ at the origin, if and only if $f$ is homogeneous. In particular,
if $X \subseteq \PP^n$ is a hypersurface given by a reduced homogeneous polynomial
$F \in \kk[x_0, \ldots, x_n]$ of degree $d$, then $X$ has multiplicity $d$ at
$p = [0: \cdots : 0: 1 : 0 : \cdots : 0]$
if and only if $F$ does not depend on $x_i$, where $i \in \{0, \ldots, n\}$
is the index of the unique non-zero coordinate of $p$.

\begin{lemma}\label{Lem:ExcDivisorMultmax}
Let $\kk$ be a field, let $n\ge 1$ and let $h\in \kk[x_0,\ldots,x_n]$ be a reduced homogeneous polynomial of degree $d\ge 1$. Let $q\in V_{\p^{n+1}}(h)(\k)$ be a $\k$-rational point of multiplicity $d$. Then, the exceptional divisor $E_q$ of the blow-up of $V_{\p^{n+1}}(h)$ at $q$ is isomorphic to $V_{\p^{n}}(h)$.
\end{lemma}
\begin{proof}
We first write $p=[0:\cdots:0:1]$ and observe that it is a point of multiplicity $d$ of $V_{\p^{n+1}}(h)$. The blow-up of $V_{\p^{n+1}}(h)$ at $p$ is then given by
\[
\Bigset{([x_0:\cdots:x_{n+1}],[y_0:\cdots:y_n])\in V_{\p^{n+1}}(h)\times V_{\p^{n}}(h)}
{\begin{array}{c}
x_iy_j=x_jy_i \\
\forall i,j\in \{0,\ldots,n\}
\end{array}},
\]
so the exceptional divisor $E_p$ of $p$ is given by $\{p\} \times V_{\p^n}(h)$ and is thus isomorphic to $V_{\p^n}(h)$. This gives the result in the case where $q=p$.

Suppose now that $q\neq p$. Applying an automorphism of $\p^{n+1}$ that fixes $p$, we may assume that $q=[0:\cdots:0:1:0]$. This implies that $h\in \k[x_0,\ldots,x_{n-1}]$. Hence, the exchange of $x_n$ and $x_{n+1}$ is an automorphism of $\p^{n+1}$ that exchanges $p$ and $q$ and preserves $V_{\p^{n+1}}(h)$. The exceptional divisors of $p$ and $q$ are then isomorphic.
\end{proof}

\begin{lemma}\label{Lem:IsoHypNNplus}
Let $\kk$ be a field, let $n\ge 1$ and let $f,g\in \kk[x_0,\ldots,x_n]$ be reduced
homogeneous polynomials of the same degree $d\ge 1$. If $V_{\p^{n+1}}(f)$ and $V_{\p^{n+1}}(g)$ are isomorphic, then $V_{\p^{n}}(f)$ and $V_{\p^{n}}(g)$ are isomorphic.
\end{lemma}
\begin{proof}
An isomorphism $V_{\p^{n+1}}(f)\iso V_{\p^{n+1}}(g)$ sends $p=[0:\cdots:0:1]$ onto a $\k$-rational point of multiplicity $d$. It then induces an isomorphism between the exceptional divisors, and thus induces an isomorphism $V_{\p^{n}}(f)\iso V_{\p^{n}}(g)$, by Lemma~\ref{Lem:ExcDivisorMultmax}.
\end{proof}

\begin{lemma}\label{LemmCurveD}
Let $\kk$ be a field, let $d\ge 4$ be an integer that is not a multiple of $\mathrm{char}(\kk)$ and let
\[f=x^{d-1}y+z^{d} \quand g=x^{d-1}y+z^{d}+\mu x^{d-2}z^2,\]
where $\mu\in \k$. Suppose that $\alpha\colon V_{\p^2}(f)\iso V_{\p^2}(g)$ is an isomorphism. Then, $\mu=0$ and $\alpha$ is equal to $[x:y:z]\mto [a^d x:y:a^{d-1} z]$ for some $a\in \k^*$.
\end{lemma}

\begin{remark}
If $d \in \{2, 3\}$ is not a multiple of $\car(\kk)$, then
there is an automorphism of $\PP^2$ that sends $V_{\PP^2}(f)$ onto $V_{\PP^2}(g)$.
\end{remark}

\begin{proof}
Expressing $y$ in terms of $x,z$, we obtain two birational morphisms
\[
\begin{aligned}
\tau\colon \p^1 & \to V_{\p^2}(f)\\
 {[u:v]} & \mto [u^d:-v^d:u^{d-1}v]
\end{aligned}\qquand
\begin{aligned}
\tau'\colon \p^1 & \to V_{\p^2}(g)\\
{[u:v]} & \mto [u^d:-v^d-\mu u^{d-2}v^2:u^{d-1}v]
\end{aligned}
\]
whose inverses are given by $[x:y:z]\mto [x:z]$. They induce isomorphisms
\[
\p^1\setminus \{[0:1]\}\iso V_{\p^2}(f)\setminus \{[0:1:0]\}\quand\p^1\setminus \{[0:1]\}\iso V_{\p^2}(g)\setminus \{[0:1:0]\}.
\]

Both $\tau$ and $\tau'$ are bijective and they send $[0:1]$ onto $[0:1:0]$, which is the unique singular point of $V_{\p^{2}}(f)$ and $V_{\p^{2}}(g)$,
respectively. Since the isomorphism $\alpha\colon V_{\p^{2}}(f)\iso V_{\p^{2}}(g)$ must fix the point $[0:1:0]$, the birational map $\hat\alpha=(\tau')^{-1}\alpha\tau$ is an automorphism of $\p^1$ that fixes the point $[0:1]$. Thus, $\hat\alpha$ is of the form $[u:v]\mto [u:av+bu]$ for some $a\in \k^*$, $b\in \k$.

Note that $s_a\colon [x:y:z]\mto [a^d x:y:a^{d-1} z]$ is an automorphism of $V_{\p^2}(f)$, that lifts to $\hat s_a=(\tau)^{-1}s_a\tau=[u:v]\mto [a u:v]$. Replacing $\alpha$ with $\alpha s_a$, we replace $\hat\alpha$ with~$\hat\alpha \hat s_a$,
and we may assume that $a=1$. It remains to see that $b=\mu=0$.
We calculate
\begin{align*}
\alpha( [x:y:z]) &= \tau' \hat\alpha \tau^{-1}( [x:y:z])\\
&= \tau' \hat\alpha([x:z]) \\
&= \tau' ([x:z+bx]) \\
&= [x^d:-(z+bx)^d-\mu x^{d-2}(z+bx)^2:x^{d-1}(z+bx)].
\end{align*}

As $\alpha$ fixes $[0:1:0]$, the pull-back by $\alpha$ of $x/y$ and $z/y$ are rational functions
\[
\frac{x^d}{-(z+bx)^d-\mu x^{d-2}(z+bx)^2} \qquand
\frac{x^{d-1}}{-(z+bx)^{d-1}-\mu x^{d-2}(z+bx)},
\]
whose restrictions to $V_{\p^2}(f)$ are regular at $[0:1:0]$. For each $j\in \{d-1,d\}$, there are thus two homogeneous polynomials $P_j,Q_j\in\k[x,y,z]$, of the same degree, such that $Q_j(0,1,0)\neq 0$ and such that \[\frac{P_j(x,y,z)}{Q_j(x,y,z)}=\frac{x^{j}}{-(z+bx)^{j}-\mu x^{d-2}(z+bx)^{j+2-d}}\] on the curve $V_{\p^2}(f)$. Using the morphism
\[
\A^1\to V_{\p^2}(f), \quad t\mto \tau([t:1])=[t^d:-1:t^{d-1}]=[t:-t^{1-d}:1],
\]
we find
\[
\frac{P_j(t^d,-1,t^{d-1})}{Q_j(t^d,-1,t^{d-1})}=\frac{t^{j}}{B_j},
\]
where $B_j = -(1+bt)^{j}-\mu t^{d-2}(1+bt)^{j+2-d} \in \kk[t]$.
As $Q_j(0,1,0)\neq 0$, the polynomial $Q_j(t^d,-1,t^{d-1})\in \k[t]$ is not divisible by~$t$. As $B_j(0) \neq 0$
(we use here $d \geq 4$), $t^j$~divides $P_j(t^d,-1,t^{d-1})$. There is thus $A_j\in \k[t]$ such that
\[
P_j(t^d,-1,t^{d-1})=A_j\cdot t^j \quad \text{and} \quad Q_j(t^d,-1,t^{d-1})=A_j\cdot B_j.
\]

We consider now the case $j=d$ and prove that $b=0$. We shall afterward prove $\mu=0$ by considering the case $j=d-1$.

As $d\ge 4$, the equality $P_d(t^d,-1,t^{d-1})=A_d\cdot t^d$ gives the existence of $\epsilon\in \k$ such that $A_d\equiv \epsilon \pmod{t^{2}}$. We moreover have $B_d\equiv -1-bdt \pmod{t^2}$. This gives
\[Q_d(t^d,-1,t^{d-1})=A_d\cdot B_d\equiv -\epsilon -\epsilon bdt \pmod{t^2}.\]
The polynomial $Q(t^d,-1,t^{d-1})$ is not divisible by $t$, so $\epsilon\neq 0$, and its coefficient of $t$ is zero, as $d\ge 4$. Since $\mathrm{char}(\k)$ does not divide $d$, this gives $b=0$.

We now use $j=d-1$. As $d\ge 4$, the equality $P_{d-1}(t^d,-1,t^{d-1})=A_{d-1}\cdot t^{d-1}$ gives the existence of $\xi,\xi'\in \k$ such that $A_{d-1}\equiv \xi+t \xi' \pmod{t^{d-1}}$. As $b=0$, we~find $B_{d-1}=B_d=-1-\mu t^{d-2}$. This gives
\[Q_{d-1}(t^d,-1,t^{d-1})=A_{d-1}\cdot B_{d-1}\equiv -\xi-t \xi'-\mu\xi t^{d-2} \pmod{t^{d-1}}.\]
As $t$ does not divide this polynomial, $\xi\neq0$. Moreover, we obtain $\xi'=0$ and $\mu=0$, as $d\ge 4$.
\end{proof}

\begin{proposition}\label{Prop.IsoHypNN3}
Let $\kk$ be a field, let $d\ge 4$ be an integer that is not a multiple of $\mathrm{char}(\kk)$ and let $f,g\in\kk[x,y,z]$ be defined by \[f=x^{d-1}y+z^{d} \text{ and } g=x^{d-1}y+z^{d}+\mu x^{d-2}z^2,\] where $\mu\in \k^*$. Then, for each $r\ge 0$, the two hypersurfaces $V_{\p^{r+2}}(f)$, $V_{\p^{r+2}}(g)$ in $\p^{r+2}$ are not isomorphic.
\end{proposition}

\begin{proof}
By Lemma~\ref{LemmCurveD}, the result is true when $r=0$. Using Lemma~\ref{Lem:IsoHypNNplus}, we can then argue by induction and obtain the result for every integer $r$.
\end{proof}

\begin{proposition}
\label{Prop.IsoHypNN4}
Let $\kk$ be a field, $d \geq 3$ be an integer and let $f,g\in\kk[x_0,x_1,y,z]$ be defined by
\[
f=x_0^{d-1}y+z^{d}, \quad g=x_0^{d-1}y+z^{d}+\mu x_0^{d-2}x_1^2,
\] where $\mu\in \k^*$. Then, for each $r\ge 0$, the two hypersurfaces $V_{\p^{r+3}}(f)$, $V_{\p^{r+3}}(g)$ in $\p^{r+3}$ are not isomorphic.
\end{proposition}

\begin{proof}
By Lemma~\ref{Lem:IsoHypNNplus}, it suffices to consider the case $r=0$ and to prove that $V_{\p^{3}}(f)$ and $V_{\p^{3}}(g)$ are not isomorphic. Looking at the derivative with respect to $y$, the singular locus of both hypersurfaces is contained in the line $\ell\subseteq \p^3$
given by $x_0 = z = 0$.
The surface $V_{\p^{3}}(f)$ has multiplicity $d$ at the point where $x_0=y=z=0$.

It remains to see that $V_{\p^{3}}(g)$ has multiplicity $< d$ at every point
and we have to check this only for points in $\ell$.
For this, write $g=z^d+ x_0^{d-2}(x_0y+\mu x_1^2)$ and observe that $V_{\PP^3}(x_0y+\mu x_1^2)$ is smooth outside $x_0=x_1=y=0$ and thus on $\ell$.
Hence, $x_0^{d-2}(x_0y+\mu x_1^2)$ has multiplicity $d-2$ or $d-1$ at any point of $\ell$, which implies that~$g$ has multiplicity $d-2$ or $d-1$ at any point of $\ell$.
\end{proof}

We may now prove Theorem~\ref{34and43}, which is a direct consequence of Propositions~\ref{Prop.Counterexample.dim3}, \ref{Prop.Counterexample.dim4}, \ref{Prop.IsoHypNN3} and \ref{Prop.IsoHypNN4}.

\begin{proof}[Proof of Theorem~\ref{34and43}]
As in the statement, we take a field $\kk$, and integers $d,n\ge 3$ such that $d$ is not a multiple of $\mathrm{char}(\kk)$ and such that $(d,n)\neq (3,3)$. We are looking for hypersurfaces $H=V_{\p^n}(f),H'=V_{\p^n}(g)\subseteq \p^n_\k$ that are not isomorphic but have isomorphic complements.

If $d\ge 4$, we may choose $f=x^{d-1}y+z^{d}$ and $g=x^{d-1}y+z^{d}+dx^{d-2}z^2$, that we
consider in $\k[x,y,z,w_1,\ldots,w_r]$ with $r=n-2\ge 1$. By Propositions~\ref{Prop.Counterexample.dim3}, the complements of $H$ and $H'$ are isomorphic, and by Proposition~\ref{Prop.IsoHypNN3}, the hypersurfaces~$H$ and $H'$ are not isomorphic.

If $n\ge 4$, we may choose $f=x_0^{d-1}y+z^{d}$ and $g=x_0^{d-1}y+z^{d}+dx_0^{d-2}x_1^2$, that we consider in
$\k[x_0,x_1,y,z,w_1,\ldots,w_r]$ with $r=n-3\ge 1$. By Propositions~\ref{Prop.Counterexample.dim4}, the complements of $H$ and $H'$ are isomorphic, and by Proposition~\ref{Prop.IsoHypNN4}, the hypersurfaces~$H$ and~$H'$ are not isomorphic.
\end{proof}

% End literal retained input author-body.tex

\begin{thebibliography}{99}
\bibitem{1973Theorie-des-topos-} M. Artin, A. Grothendieck \& J.-L. Verdier - Théorie des topos et cohomologie étale des schémas. Tome 3 , Lect. Notes in Math. , vol. 305 , Springer-Verlag, Berlin-New York, 1972–1973, Séminaire de Géométrie Algébrique du Bois–Marie 1963–64 (SGA 4)
\bibitem{AtMa1969Introduction-to-co} M. F. Atiyah \& I. G. Macdonald - Introduction to commutative algebra , Addison-Wesley Publishing Co., Reading, Mass.-London-Don Mills, Ont., 1969
\bibitem{Ar1962Some-numerical-cri} M. Artin - “Some numerical criteria for contractability of curves on algebraic surfaces” , Amer. J. Math. 84 (1962), p. 485-496
\bibitem{BlFaTe2023Connected-Algebrai} J. Blanc, A. Fanelli \& R. Terpereau - “Connected algebraic groups acting on three-dimensional Mori fibrations” , Internat. Math. Res. Notices (2023) no. 2, p. 1572-1689
\bibitem{BlancC} J. Blanc - “The correspondence between a plane curve and its complement” , J. reine angew. Math. 633 (2009), p. 1-10
\bibitem{Bo2018The-class-of-the-a} L. A. Borisov - “The class of the affine line is a zero divisor in the Grothendieck ring” , J. Algebraic Geom. 27 (2018) no. 2, p. 203-209
\bibitem{CheDub} I. Cheltsov, A. Dubouloz \& J. Park - “Super-rigid affine Fano varieties” , Compositio Math. 154 (2018) no. 11, p. 2462-2484
\bibitem{Costa} P. Costa - “New distinct curves having the same complement in the projective plane” , Math. Z. 271 (2012) no. 3-4, p. 1185-1191
\bibitem{Cr2004On-the-AK-invarian} A. J. Crachiola - On the AK invariant of certain domains , Ph. D. Thesis, Wayne State University, Ann Arbor, MI, 2004, 84 pages, ProQuest LLC
\bibitem{2001Correspondance-Gro} - Correspondance Grothendieck-Serre (P. Colmez \& J.-P. Serre, eds.) , Documents Mathématiques , vol. 2 , Société Mathématique de France, Paris, 2001
\bibitem{DuKi2015Log-uniruled-affin} A. Dubouloz \& T. Kishimoto - “Log-uniruled affine varieties without cylinder-like open subsets” , Bull. Soc. math. France 143 (2015) no. 2, p. 383-401
\bibitem{Gr1961Elements-de-geomet-III} A. Grothendieck - “Éléments de géométrie algébrique. III. Étude cohomologique des faisceaux cohérents. I” , Publ. Math. Inst. Hautes Études Sci. (1961) no. 11, p. 167
\bibitem{2005Affine-algebraic-g} - Affine algebraic geometry (J. Gutierrez, V. Shpilrain \& J.-T. Yu, eds.) , Contemp. Math. , vol. 369 , American Mathematical Society, Providence, RI, 2005, Papers from the Special Session at the 1st Joint AMS-RSME Meeting held in Seville, June 18–21, 2003
\bibitem{Ha1977Algebraic-geometry} R. Hartshorne - Algebraic geometry , Graduate Texts in Math , vol. 52 , Springer-Verlag, New York-Heidelberg, 1977
\bibitem{Hemmig} M. Hemmig - “Isomorphisms between complements of projective plane curves” , Épijournal de Géom. Alg. 3 (2019), article ID 16, 44 pages
\bibitem{KuSh2018Grothendieck-ring-} A. Kuznetsov \& E. Shinder - “Grothendieck ring of varieties, D- and L-equivalence, and families of quadrics” , Selecta Math. (N.S.) 24 (2018) no. 4, p. 3475-3500
\bibitem{KoSmCo2004Rational-and-nearl} J. Kollár, K. E. Smith \& A. Corti - Rational and nearly rational varieties , Cambridge Studies in Advanced Math. , vol. 92 , Cambridge University Press, Cambridge, 2004
\bibitem{Li2005Lectures-on-topolo} A. Libgober - “Lectures on topology of complements and fundamental groups” , 2005
\bibitem{LaLu2003Motivic-measures-a} M. Larsen \& V. A. Lunts - “Motivic measures and stable birational geometry” , Moscow Math. J. 3 (2003) no. 1, p. 85-95
\bibitem{LePaSc2011On-the-classificat} W. Lee, E. Park \& P. Schenzel - “On the classification of non-normal cubic hypersurfaces” , J. Pure Appl. Algebra 215 (2011) no. 8, p. 2034-2042
\bibitem{LiSe2010The-Grothendieck-r} Q. Liu \& J. Sebag - “The Grothendieck ring of varieties and piecewise isomorphisms” , Math. Z. 265 (2010) no. 2, p. 321-342
\bibitem{LaSe2012Birational-self-ma} S. Lamy \& J. Sebag - “Birational self-maps and piecewise algebraic geometry” , J. Math. Sci. Univ. Tokyo 19 (2012) no. 3, p. 325-357
\bibitem{Ma2016The-class-of-the-a} N. Martin - “The class of the affine line is a zero divisor in the Grothendieck ring: an improvement” , Comptes Rendus Mathématique 354 (2016) no. 9, p. 936-939
\bibitem{Ma1986Commutative-ring-t} H. Matsumura - Commutative ring theory , Cambridge Studies in Advanced Math. , vol. 8 , Cambridge University Press, Cambridge, 1986
\bibitem{Mi1980Etale-cohomology} J. S. Milne - Étale cohomology , Princeton Math. Series , vol. 33 , Princeton University Press, Princeton, NJ, 1980
\bibitem{Mi2013Lectures-on-Etale-} J. S. Milne - “Lectures on Etale Cohomology (v2.21)” , 2013, p. 202, Available at www.jmilne.org/math/
\bibitem{MoPo2021Isomorphisms-betwe} L. Moser-Jauslin \& P.-M. Poloni - “Isomorphisms between cylinders over Danielewski surfaces” , Beitr. Algebra Geom. 62 (2021) no. 4, p. 755-771
\bibitem{Pf1995Quadratic-forms-wi} A. Pfister - Quadratic forms with applications to algebraic geometry and topology , London Math. Soc. Lect. Note Series , vol. 217 , Cambridge University Press, Cambridge, 1995
\bibitem{PoloniClassifications} P.-M. Poloni - “Classification(s) of Danielewski hypersurfaces” , Transform. Groups 16 (2011) no. 2, p. 579-597
\bibitem{Sakamaki:2010vx} Y. Sakamaki - “Automorphism groups on normal singular cubic surfaces with no parameters” , Trans. Amer. Math. Soc. 362 (2010) no. 5, p. 2641-2666
\bibitem{Yoshihara} H. Yoshihara - “On open algebraic surfaces ${\bf P}^{2}-C$” , Math. Ann. 268 (1984) no. 1, p. 43-57
\end{thebibliography}
\end{document}
